\section{Introduction and claim ledger}
\label{sec:introduction}

The factor-scale coordinate introduced in TPC--9 and TPC--10 records a
factorization $r=mn$ by
\[
 u_{m,n}=\log(n/m).
\]
At frequency zero, summing over all divisors gives the exact collapse
\[
 \sum_{n\mid r}(\Lambda(n)-1)=\log r-\tau(r),
\]
which leads to a fixed-shift Titchmarsh evaluation
\citep{AssingBlomerLi2021,WangTPC9}.  TPC--10 proved coarse two-sided
factor localization and arbitrary fixed smooth localization on the Type~I
side.  It also showed that finitely many Mellin coordinates do not stably
determine a local factor-scale window, and proposed a continuum-band mean
square as the next analytic target \citep{WangTPC10}.

The present paper treats that target in two different quantifier regimes.
Let $X\geq2$, $h\in\Z$, and $D\geq1$.  Fix throughout
\begin{equation}
 u(n):=\Lambda(n)-1,
 \qquad W\in C_c^1((0,\infty);\R),
 \qquad \supp W\subset[\alpha,\beta]
 \label{eq:intro-basic-data}
\end{equation}
with $0<\alpha<\beta<\infty$.  We extend $\Lambda(q)$ by zero for
$q\leq0$.  For any integrable function $F$ on $(0,\infty)$, put
\begin{equation}
 I_0(F):=\int_0^\infty F(v)\,dv,
 \qquad
 I_2(F):=\int_0^\infty |F(v)|^2\,dv.
 \label{eq:moment-functionals}
\end{equation}
All asymptotic statements are taken as $X\to\infty$; parameters declared
fixed in a theorem are held fixed, while $D,H$, and $T$ may vary in the
displayed ranges.
Define
\begin{align}
 \cT_{h,D}(X;t)
 &:={}
 \sum_{\substack{n>D\\m\geq1}}
 u(n)\Lambda(mn+h)W(mn/X)(n/m)^{it},
 \label{eq:intro-raw-tail}\\
 \cM_D(X;t)
 &:={}
 \sum_{\substack{n>D\\m\geq1}}
 u(n)W(mn/X)(n/m)^{it},
 \label{eq:intro-source-model}\\
 \cR_{h,D}(X;t)
 &:={}\cT_{h,D}(X;t)-\cM_D(X;t).
 \label{eq:intro-centered-tail}
\end{align}
Thus $\cR_{h,D}$ is obtained by replacing the target factor
$\Lambda(mn+h)$ by $\Lambda(mn+h)-1$.  The subtraction is algebraic for
every fixed $h$.  It becomes a genuine first-moment model only after the
long shift average proved in \cref{prop:first-moment-model}.

\subsection{What is proved}

There are four main outputs.

\begin{enumerate}[label=\textup{(\roman*)}]
 \item Every factor pair admits a unique representation
 \[
  n=ak,\qquad m=bk,\qquad (a,b)=1.
 \]
 Equal ratios are thereby compressed into one primitive ray $(a,b)$.
 Distinct ray frequencies $\log(a/b)$ have spacing $\gg_W X^{-1}$.
 Applying the continuous mean-value theorem of
 \citet{MontgomeryVaughan1974} gives an exact fixed-shift diagonal law up
 to an $O(X)$ resolution error.

 \item The ray energy has the unconditional fixed-shift bound
 \[
  \cE_{h,D}(X)\ll_{C,W}X(\log X)^4
  \qquad (|h|\leq X^C).
 \]
 Hence a band of length $T\gg(\log X)^4$ gives $o(X)$ cancellation in
 root-mean-square per frequency, for every prescribed shift.  This does
 not evaluate any specified Mellin frequency.

 \item After the error is squared, a smooth average over shifts
 $h\asymp H$, $H\geq X$, permits a full arithmetic evaluation of the ray
 energy.  The $k=\ell$ target diagonal has size $H\log H$.  Same-ray
 terms $k\ne\ell$ require prime-pair upper bounds, but not prime-pair
 asymptotics: a Selberg upper sieve costs only $H\log\log X$, and a
 source-ray ledger sums their coefficients to $O_W(X\log^2X)$.  The
 diagonal therefore wins by one logarithm.

 \item A compact-Fourier-support kernel resolves the factor rays exactly
 at the natural threshold $T\asymp X$.  Below that threshold, the
 surviving interactions are indexed by the integer determinant
 $ad-bc$.  At the opposite endpoint, a fixed smooth scale window can
 select $m=1$ exactly on every product band of multiplicative width less
 than four.  This gives a quantitative $o(X^2)$ strength threshold for
 any fixed-shift continuum-band model that could imply the corresponding
 Hardy--Littlewood functional.
\end{enumerate}

The generalized Hilbert inequality is the only input in the fixed-shift
spectral theorem.  The averaged-shift energy theorem additionally uses
the prime number theorem and a standard two-linear-forms Selberg upper
sieve.  It does not import an averaged Hardy--Littlewood asymptotic.
For comparison, Theorem~1.3(i) of
\citet{MatomakiRadziwillTao2019} evaluates two-point prime correlations for
almost all shifts in the substantially shorter range
$X^{8/33+\eps}\leq H\leq X^{1-\eps}$; TPC--7 already transferred that
result to a different complete defect \citep{WangTPC7}.  The later
higher-uniformity theorem of
\citet[Theorem~1.5]{MatomakiRadziwillShaoTaoTeravainen2026} treats higher
orders while retaining the earlier two-point input.  Our theorem instead
evaluates a continuum of Mellin frequencies after a longer second-moment
shift average.

\subsection{What is not proved}

The order of averaging is essential.  We first form
$|\cR_{h,D}(X;t)|^2$ and only then average $h$.  A coherent average of
$\cR_{h,D}$ before squaring is much easier and says nothing about a
typical shift.  Likewise, an average over Mellin frequencies does not
select $t=0$ or any other prescribed frequency.

On one primitive ray, the target arguments are
\[
 abk^2+h.
\]
No amount of separation between distinct values of $a/b$ decorrelates
different $k$ on the same ray.  In particular, the ray $(1,1)$ already
contains values of $k^2+h$, while the factor endpoint $b=1$, $k=1$
contains the fixed-shift prime-pair functional.  We neither estimate these
objects at a prescribed $h$ nor assert that such an estimate is
impossible.  No twin-prime lower bound, fixed-shift Hardy--Littlewood
asymptotic, new Type~II estimate, new distribution level, or parity-barrier
breakthrough follows.

For later reference, the logical ledger is displayed in
\cref{tab:intro-ledger}.

\begin{table}[ht]
\centering
\small
\begin{tabular}{p{0.27\textwidth}p{0.33\textwidth}p{0.31\textwidth}}
\toprule
Input & Proved output & Not obtained\\
\midrule
Reduced factor ratios + Montgomery--Vaughan
 & Fixed-$h$ continuum mean square and exact ray resolution
 & Arithmetic evaluation of one fixed-$h$ ray energy\\
Long smooth shift average + PNT + upper sieve
 & Explicit averaged ray-energy and band asymptotics
 & Any prescribed individual shift\\
Compact-Fourier kernel
 & Exact removal of distinct rays at $T\asymp X$
 & Separation of different $k$ on one ray\\
Endpoint scale selector
 & Exact $m=1$ reduction and an $L^2$ strength threshold
 & The Hardy--Littlewood value of that endpoint\\
\bottomrule
\end{tabular}
\caption{Every implication used in the paper is one-way.}
\label{tab:intro-ledger}
\end{table}
